% !TEX program = xelatex
% 慧学外语阅读自动答题 — 零基础安装与使用指南
% 适用于 Windows 和 macOS 平台的完全新手教程
\documentclass[12pt,a4paper]{ctexart}

% ========== 页面设置 ==========
\usepackage[top=2.5cm,bottom=2.5cm,left=2.8cm,right=2.8cm]{geometry}
\usepackage[hidelinks]{hyperref}
\usepackage{enumitem}
\usepackage{tcolorbox}
\usepackage{fancyhdr}
\usepackage{listings}
\usepackage{xcolor}
\usepackage{fontawesome5}
\usepackage{titlesec}
\usepackage{cleveref}

% ========== 页眉高度 ==========
\setlength{\headheight}{14pt}
\addtolength{\topmargin}{-2pt}

% ========== 颜色定义 ==========
\definecolor{primary}{HTML}{2563EB}
\definecolor{success}{HTML}{16A34A}
\definecolor{warning}{HTML}{EA580C}
\definecolor{danger}{HTML}{DC2626}
\definecolor{graybg}{HTML}{F3F4F6}
\definecolor{codebg}{HTML}{1E293B}

% ========== 代码块样式 ==========
\lstset{
  basicstyle=\ttfamily\footnotesize\color{white},
  backgroundcolor=\color{codebg},
  frame=single,
  rulecolor=\color{codebg},
  framerule=2pt,
  framexleftmargin=6pt,
  framexrightmargin=6pt,
  xleftmargin=8pt,
  xrightmargin=8pt,
  aboveskip=10pt,
  belowskip=10pt,
  breaklines=true,
  breakatwhitespace=false,
  showstringspaces=false,
  columns=flexible,
  keepspaces=true,
  language=bash,
  keywordstyle=\color{cyan},
  commentstyle=\color{green!40!white},
  stringstyle=\color{yellow!60!white},
  tabsize=2,
}

% ========== 提示框 ==========
\newtcolorbox{tipbox}[1][]{
  colback=graybg,
  colframe=primary,
  leftrule=4mm,
  arc=3mm,
  boxrule=0.5mm,
  title=#1,
  fonttitle=\bfseries\color{primary},
  before skip=12pt,
  after skip=12pt,
}

\newtcolorbox{warnbox}[1][]{
  colback=red!5,
  colframe=warning,
  leftrule=4mm,
  arc=3mm,
  boxrule=0.5mm,
  title=#1,
  fonttitle=\bfseries\color{warning},
  before skip=12pt,
  after skip=12pt,
}

% ========== 页眉页脚 ==========
\pagestyle{fancy}
\fancyhf{}
\fancyhead[L]{\small\color{gray}慧学外语阅读自动答题 —— 零基础指南}
\fancyhead[R]{\small\color{gray}\thepage}
\renewcommand{\headrulewidth}{0.4pt}

% ========== 标题格式 ==========
\titleformat{\section}{\Large\bfseries\color{primary}}{}{0em}{}[\titlerule]
\titleformat{\subsection}{\large\bfseries}{}{0em}{}
\titleformat{\subsubsection}{\normalsize\bfseries}{}{0em}{}

% ========== 键盘快捷键 ==========
\newcommand{\kb}[1]{\texttt{#1}}
\newcommand{\winR}{\faWindows{} + R}
\newcommand{\cmdSpace}{\faApple{}Command + 空格}

% ========== 文档信息 ==========
\title{{\Huge\bfseries 慧学外语阅读}\\[2mm]{\Large 零基础安装与使用指南}}
\author{%
  \small 适用平台：Windows 10/11 \quad$\bullet$\quad macOS 12+ \\[2mm]
  \small 最后更新：\today
}
\date{}

\begin{document}

\maketitle
\thispagestyle{empty}

\vspace{10mm}

\begin{center}
  \fbox{
    \begin{minipage}{0.85\textwidth}
      \small\textbf{\faLightbulb{} 这份教程适合谁？}
      \begin{itemize}[nosep]
        \item 完全不懂编程的同学
        \item 没装过 Python，不知道``环境变量''是什么
        \item 只想一键刷完英语阅读，不想折腾技术细节
      \end{itemize}
      \textbf{你需要准备：}
      \begin{itemize}[nosep]
        \item 一台能上网的电脑（Windows 或 Mac）
        \item ZJU 统一认证账号（学号 + 密码）
        \item 大约 30 分钟的时间
        \item 耐心——按照步骤操作，每一步都有说明
      \end{itemize}
  \end{minipage}}
\end{center}

\vfill
\newpage

% ============================================================
\tableofcontents
\newpage

% ============================================================
\section{在开始之前}

\subsection{这个工具是什么？}

简单说：一个自动帮你完成``慧学外语''平台上所有英语阅读题的程序。你不需要手动点开每篇文章、阅读、答题——它全部自动完成。

\subsection{它怎么工作的？}

\begin{enumerate}[nosep]
  \item 你运行一个命令，它打开 Edge 浏览器
  \item 自动登录你的 ZJU 账号
  \item 自动浏览 11 个阅读主题（约 291 篇文章）
  \item 每篇文章自动提取内容，配合 AI 答题
  \item 答完自动提交，继续下一篇
\end{enumerate}

\bigskip

\begin{warnbox}[\faExclamationTriangle{} 重要提醒]
  \begin{itemize}[nosep]
    \item 运行期间\textbf{不要关闭}自动打开的浏览器窗口
    \item 运行期间\textbf{不要操作}浏览器（不要点击、不要切换标签页）
    \item 仅供学习用途，请合理使用
    \item 正常刷完所有文章约需 1--2 小时
  \end{itemize}
\end{warnbox}

% ============================================================
\newpage
\section{Windows 平台完整指南}

\subsection{第一步：确认你的 Windows 版本}

\begin{enumerate}[nosep]
  \item 按键盘上的 \winR
  \item 输入 \texttt{winver} 然后按回车
  \item 确认版本是 Windows 10 或 Windows 11
\end{enumerate}

\subsection{第二步：下载项目文件}

\begin{enumerate}
  \item 打开浏览器，访问：
    \begin{lstlisting}[language=]
https://github.com/pirate-608/huixuewaiyu-readingpart
\end{lstlisting}
  \item 点击页面右侧绿色的``\textbf{克隆/下载}''按钮
  \item 选择``\textbf{下载 ZIP}''
  \item 下载完成后，找到 \texttt{huixuewaiyu-readingpart.zip} 文件
  \item \textbf{右键} → ``\textbf{全部解压缩}'' →
    选择一个容易找到的位置（比如桌面），点击``\textbf{解压缩}''
\end{enumerate}

\begin{tipbox}[\faFolderOpen{} 解压到哪？]
  建议解压到桌面，解压后你会看到一个名为 \texttt{huixuewaiyu-readingpart}
  的文件夹。后面所有操作都在这个文件夹里进行。
\end{tipbox}

\subsection{第三步：安装 Python}

\begin{enumerate}
  \item 打开浏览器，访问：
    \begin{lstlisting}[language=]
https://www.python.org/downloads/
\end{lstlisting}
  \item 页面会自动识别你的系统，显示``\textbf{Download Python 3.12.x}''（版本号可能略有不同）
  \item 点击黄色下载按钮，下载安装程序
  \item \textbf{双击}下载的 \texttt{python-3.12.x-amd64.exe}
  \item \textbf{关键步骤}：在安装窗口的\textbf{最底部}，\textbf{勾选}
    ``\texttt{Add Python to PATH}''（把 Python 添加到环境变量）
  \item 然后点击``\textbf{Install Now}''（立即安装）
  \item 等待安装完成，点击``\textbf{Close}''
\end{enumerate}

\begin{warnbox}[\faExclamationTriangle{} 最容易出错的步骤]
  \textbf{一定}要勾选 ``\texttt{Add Python to
  PATH}''！如果不勾选，后面所有命令都会报``找不到 Python''的错误。\\
  如果忘记勾选了：重新运行安装程序 → 选择``\textbf{Modify}''（修改）→ 勾选
  ``\texttt{Add Python to PATH}'' → 安装。
\end{warnbox}

\subsection{第四步：验证 Python 安装}

\begin{enumerate}
  \item 按 \winR，输入 \texttt{cmd}，按回车（打开命令提示符）
  \item 在黑色窗口中输入以下命令，然后按回车：
    \begin{lstlisting}[language=]
python --version
\end{lstlisting}
  \item 如果看到类似 ``\texttt{Python 3.12.x}'' 的输出，说明安装成功
  \item 如果看到 ``\texttt{'python' 不是内部或外部命令}''，说明第三步的
    ``\texttt{Add Python to PATH}'' 没有勾选，请回到第三步重新安装
\end{enumerate}

\subsection{第五步：运行安装脚本}

\begin{enumerate}
  \item 打开你解压的 \texttt{huixuewaiyu-readingpart} 文件夹
  \item 在文件夹的\textbf{空白处}，按住 \textbf{Shift} 键，然后\textbf{右键单击}
  \item 在弹出菜单中选择``\textbf{在此处打开 PowerShell 窗口}''（或``在此处打开终端''）
  \item 在打开的蓝色（或黑色）窗口中，输入以下命令，然后按回车：
    \begin{lstlisting}[language=]
powershell -ExecutionPolicy Bypass -File install.ps1
\end{lstlisting}
  \item 脚本会自动完成以下操作：
    \begin{itemize}[nosep]
      \item 检查 Python \faCheck{}
      \item 创建独立虚拟环境 \faCheck{}
      \item 安装所有依赖 \faCheck{}
      \item 下载 Chromium 浏览器驱动 \faCheck{}
      \item 复制技能文件到 Claude Code 目录 \faCheck{}
    \end{itemize}
  \item 中途会提示输入\textbf{学号}和\textbf{CAS 密码}——这是登录慧学外语用的，仅存储在本地
\end{enumerate}

\begin{tipbox}[\faLock{} 凭据安全]
  你的学号和密码保存在：
  \begin{center}
    \texttt{C:\textbackslash Users\textbackslash
      你的用户名\textbackslash .claude\textbackslash
    skills\textbackslash huixuewaiyu-readingpart\textbackslash .env}
  \end{center}
  不会上传到任何地方。
\end{tipbox}

\subsection{第六步：开始刷题}

安装完成后，在 \textbf{Claude Code} 中输入：

\begin{lstlisting}[language=]
/huixuewaiyu-readingpart
\end{lstlisting}

或者直接说``\textbf{慧学外语刷题}''。

程序会自动开始工作，浏览器窗口会自动打开，你会看到它自动：
\begin{itemize}[nosep]
  \item 登录 CAS
  \item 浏览各个阅读主题
  \item 一篇篇点开文章、答题、提交
\end{itemize}

\bigskip

\begin{warnbox}[\faShield*{} 验证码]
  连续刷约 10 篇文章后，平台会弹出 4
  位数字验证码，程序会自动识别。极少数情况下识别失败，需要你手动输入——\textbf{注意看浏览器窗口}，程序会等你 5 分钟。
\end{warnbox}

\begin{tipbox}[\faChartBar{} 断点续传]
  如果中途想停下来，程序每完成 50 篇会暂停询问。你也可以直接关闭窗口。下次再运行会自动从上次的位置继续，不会重复答题。
\end{tipbox}

% ============================================================
\newpage
\section{macOS 平台完整指南}

\subsection{第一步：确认你的 macOS 版本}

\begin{enumerate}[nosep]
  \item 点击屏幕左上角 \faApple{} 图标
  \item 选择``\textbf{关于本机}''
  \item 确认版本是 macOS 12 (Monterey) 或更高
\end{enumerate}

\subsection{第二步：下载项目文件}

\begin{enumerate}
  \item 打开 Safari 浏览器，访问：
    \begin{lstlisting}[language=]
https://github.com/pirate-608/huixuewaiyu-readingpart
\end{lstlisting}
  \item 点击页面右侧绿色的``\textbf{克隆/下载}''按钮
  \item 选择``\textbf{下载 ZIP}''
  \item 下载完成后，Safari 会自动解压到``\textbf{下载}''文件夹
  \item 打开 Finder，进入``\textbf{下载}''文件夹，找到
    \texttt{huixuewaiyu-readingpart} 文件夹
  \item 把它\textbf{拖到桌面}（方便后续操作）
\end{enumerate}

\subsection{第三步：安装 Python}

macOS 通常自带 Python，但版本可能较旧。我们安装最新版。

\begin{enumerate}
  \item 打开 Safari，访问：
    \begin{lstlisting}[language=]
https://www.python.org/downloads/
\end{lstlisting}
  \item 页面会自动识别你在用 Mac，显示``\textbf{Download Python 3.12.x for macOS}''
  \item 点击黄色下载按钮，下载 \texttt{python-3.12.x-macos11.pkg}
  \item \textbf{双击}下载的 \texttt{.pkg} 文件
  \item 按照安装向导操作：点击``\textbf{继续}'' → ``\textbf{继续}'' →
    ``\textbf{同意}'' → ``\textbf{安装}''
  \item 输入你的 Mac 登录密码，点击``\textbf{安装软件}''
  \item 安装完成后点击``\textbf{关闭}''
\end{enumerate}

\subsection{第四步：验证 Python 安装}

\begin{enumerate}
  \item 按 \cmdSpace{}，输入 \texttt{终端}（Terminal），按回车
  \item 在打开的白色窗口中输入以下命令，按回车：
    \begin{lstlisting}[language=]
python3 --version
\end{lstlisting}
  \item 如果看到 ``\texttt{Python 3.12.x}''，说明安装成功
\end{enumerate}

\subsection{第五步：安装 Edge 浏览器}

\begin{enumerate}
  \item 打开 Safari，访问：
    \begin{lstlisting}[language=]
https://www.microsoft.com/edge/download
\end{lstlisting}
  \item 点击``\textbf{下载 macOS 版}''
  \item \textbf{双击}下载的 \texttt{.pkg} 文件
  \item 按照安装向导完成安装
\end{enumerate}

\begin{tipbox}[\faQuestionCircle{} 为什么需要 Edge？]
  这个工具使用 Playwright 自动化框架，它通过 Edge 浏览器来控制慧学外语网页。Chrome
  也可以，但 Edge 兼容性更好（尤其是在国内网络环境下）。
\end{tipbox}

\subsection{第六步：运行安装脚本}

\begin{enumerate}
  \item 打开\textbf{终端}（\cmdSpace{} → 输入 \texttt{终端}）
  \item 输入以下命令进入项目文件夹（假设你放在了桌面）：
    \begin{lstlisting}[language=]
cd ~/Desktop/huixuewaiyu-readingpart
\end{lstlisting}
  \item 运行安装脚本：
    \begin{lstlisting}[language=]
bash install.sh
\end{lstlisting}
  \item 脚本会自动完成：
    \begin{itemize}[nosep]
      \item 检查 Python \faCheck{}
      \item 创建独立虚拟环境 \faCheck{}
      \item 安装所有依赖 \faCheck{}
      \item 下载 Chromium 浏览器驱动 \faCheck{}
      \item 复制技能文件到 Claude Code 目录 \faCheck{}
    \end{itemize}
  \item 中途会提示输入\textbf{学号}和\textbf{CAS 密码}
\end{enumerate}

\begin{tipbox}[\faLock{} 凭据安全]
  你的学号和密码保存在
  \texttt{\textasciitilde/.claude/skills/huixuewaiyu-readingpart/.env}，仅存储在本地。
\end{tipbox}

\subsection{第七步：开始刷题}

安装完成后，在 \textbf{Claude Code} 中输入：

\begin{lstlisting}[language=]
/huixuewaiyu-readingpart
\end{lstlisting}

或者直接说``\textbf{慧学外语刷题}''。

% ============================================================
\newpage
\section{常见问题与故障排除}

\subsection{Windows 常见问题}

\begin{tipbox}[Q: PowerShell 提示``无法加载文件''或``禁止执行脚本'']
  \textbf{A:} 这是 PowerShell 的安全策略。用管理员身份打开 PowerShell，运行：
    \begin{lstlisting}[language=]
Set-ExecutionPolicy -ExecutionPolicy RemoteSigned -Scope CurrentUser
\end{lstlisting}
  然后输入 \texttt{Y} 确认。关闭窗口重新打开即可。
\end{tipbox}

\begin{tipbox}[Q: 提示 ``python 不是内部或外部命令'']
  \textbf{A:} Python 没有添加到 PATH。解决方案：
  \begin{enumerate}[nosep]
    \item 搜索栏输入``\textbf{应用和功能}''，找到 Python，点击``\textbf{修改}''
    \item 在安装界面勾选 ``\texttt{Add Python to PATH}''
    \item 或者重新安装 Python，这次一定勾选那个选项
  \end{enumerate}
\end{tipbox}

\begin{tipbox}[Q: pip 安装报错，提示 ``opencv-python'' 相关]
  \textbf{A:} 你安装的 Python 版本可能太新。请确认安装的是 Python 3.12。
\end{tipbox}

\begin{tipbox}[Q: Edge 浏览器没有自动打开]
  \textbf{A:} 请确认已经安装了 Microsoft Edge 浏览器（Windows 10/11
  自带）。如果没有，打开 \url{https://www.microsoft.com/edge/download} 下载并安装。
\end{tipbox}

\subsection{macOS 常见问题}

\begin{tipbox}[Q: 提示 ``command not found: python3'']
  \textbf{A:} Python 没有正确安装。重新运行
  \texttt{python-3.12.x-macos11.pkg} 安装程序。
\end{tipbox}

\begin{tipbox}[Q: 提示 ``SSL 证书错误'' 或 ``CERTIFICATE\_VERIFY\_FAILED'']
  \textbf{A:} 在终端运行以下命令安装证书：
    \begin{lstlisting}[language=]
/Applications/Python\ 3.12/Install\ Certificates.command
\end{lstlisting}
\end{tipbox}

\begin{tipbox}[Q: Edge 浏览器没有自动打开]
  \textbf{A:} 请确认已经下载安装 Microsoft Edge（见 macOS 第五步）。
\end{tipbox}

\subsection{通用问题}

\begin{tipbox}[Q: 程序运行到一半停了，浏览器没反应]
  \textbf{A:}
  \begin{enumerate}[nosep]
    \item 检查浏览器窗口——可能需要你手动输入验证码
    \item 如果弹出``\textbf{验证码}''窗口，程序在等你输入（最多等 5 分钟）
    \item 不要手动点击浏览器的任何按钮
  \end{enumerate}
\end{tipbox}

\begin{tipbox}[Q: 提示 ``.env 文件不存在'' 或 ``CAS\_USERNAME 未设置'']
  \textbf{A:} 安装脚本中途会询问学号和密码。如果跳过了，手动创建 \texttt{.env} 文件：
  \begin{enumerate}[nosep]
    \item 在项目文件夹中找到 \texttt{.env.example}
    \item 复制一份，重命名为 \texttt{.env}
    \item 用记事本（或文本编辑）打开，填入你的学号和密码
  \end{enumerate}
\end{tipbox}

\begin{tipbox}[Q: 我已经刷了一些，想重新从头开始刷]
  \textbf{A:} 删除断点续传文件：
  \begin{itemize}[nosep]
    \item Windows：删除 \texttt{C:\textbackslash
      tmp\textbackslash elang\_checkpoint.json}
    \item macOS：删除 \texttt{/tmp/elang\_checkpoint.json}
  \end{itemize}
\end{tipbox}

\begin{tipbox}[Q: 想只刷某一个主题，不想全部刷]
  \textbf{A:} 在 Claude Code
  中说``\textbf{慧学外语刷题}''，然后告诉它你想刷的主题编号（1--11）。或者直接在终端运行：
    \begin{lstlisting}[language={},breaklines=true]
# Windows (PowerShell)
& "$env:USERPROFILE\.claude\skills\huixuewaiyu-readingpart\.venv\Scripts\python.exe" "$env:USERPROFILE\.claude\skills\huixuewaiyu-readingpart\scripts\elang_reader.py" batch "https://elang.zju.edu.cn/#/read/learn?subject_id=6"
\end{lstlisting}
    \begin{lstlisting}[language={},breaklines=true]
# macOS (终端)
~/.claude/skills/huixuewaiyu-readingpart/.venv/bin/python ~/.claude/skills/huixuewaiyu-readingpart/scripts/elang_reader.py batch "https://elang.zju.edu.cn/#/read/learn?subject_id=6"
\end{lstlisting}
  主题编号：1=道路与交通, 2=历史与文化, 3=文学与艺术, 4=职业与发展, 5=运动与娱乐,
  6=学习与教育, 7=商业与经济, 8=科技与创新, 9=社会与政治, 10=自然与农业, 11=家庭与生活。
\end{tipbox}

% ============================================================
\newpage
\section{主题编号速查表}

\begin{center}
  \begin{tabular}{|c|l|c|}
    \hline
    \textbf{编号} & \textbf{主题名称} & \textbf{文章数} \\
    \hline
    1 & 道路与交通 & 3 \\
    2 & 历史与文化 & 22 \\
    3 & 文学与艺术 & 12 \\
    4 & 职业与发展 & 18 \\
    5 & 运动与娱乐 & 6 \\
    6 & 学习与教育 & 59 \\
    7 & 商业与经济 & 26 \\
    8 & 科技与创新 & 38 \\
    9 & 社会与政治 & 36 \\
    10 & 自然与农业 & 22 \\
    11 & 家庭与生活 & 49 \\
    \hline
    \multicolumn{2}{|r|}{\textbf{合计}} & \textbf{\textasciitilde291} \\
    \hline
  \end{tabular}
\end{center}

\vspace{10mm}

\begin{center}
  \large\color{primary}\textbf{祝刷题顺利！\faSmile{}}
\end{center}

\end{document}
